% elegantOverview.tex
%
% A standalone LaTeX rendering of the "Overview of elegant and SDDS" page,
% https://ops.aps.anl.gov/elegant.html (Michael Borland, 2010-12-20).
%
% Build with:  pdflatex elegantOverview.tex
%
\documentclass[11pt]{article}

\usepackage[margin=1in]{geometry}
\usepackage[T1]{fontenc}
\usepackage{enumitem}
\usepackage[hidelinks]{hyperref}

\newcommand{\elegant}{\texttt{elegant}}
\newcommand{\sdds}{SDDS}

\title{Overview of \elegant{} and \sdds{}}
\author{Michael Borland}
\date{2010-12-20}

\begin{document}
\maketitle

\begin{abstract}
This document briefly summarizes the capabilities of \elegant{} and the
\sdds{} toolkit for simulating storage rings, linacs, and beam-transport
systems, in both the serial (\elegant{}) and parallel (\texttt{Pelegant})
versions. It is intended for potential new users who want a general sense of
what the codes can do without first reading the manual. Detailed examples are
not given here; instead, the reader is referred to the set of example files.
This is a LaTeX rendering of the online overview page,
\url{https://ops.aps.anl.gov/elegant.html}.
\end{abstract}

\section{Introduction}

\elegant{} (ELEctron Generation ANd Tracking) is an accelerator-simulation
program that has grown to address a wide variety of single-particle and
collective-dynamics problems. It is open source: both the source code and
ready-to-run executables are available. New users are encouraged to obtain the
example files and to join the user forum. At a minimum, users should subscribe
to the forum's ``Bugs'' topic in order to receive notification of bugs and
fixes.

The capabilities of \elegant{} and its associated tools fall into three broad
categories, described in the sections that follow:
\begin{itemize}[nosep]
  \item \textbf{Simulation capabilities}, which correspond to commands in the
        main input file.
  \item \textbf{Physics capabilities}, which correspond to the elements used in
        the lattice-definition file.
  \item \textbf{Cooperative capabilities}, which are provided by other programs
        that work together with \elegant{} through \sdds{} files.
\end{itemize}

\paragraph{Where to get more.} The following resources are available online:
\begin{itemize}[nosep]
  \item Example files:
        \url{https://ops.aps.anl.gov/cgi-bin/oagLog4.cgi?name=elegantExamples.tar.gz}
  \item Source code and executables:
        \url{https://www.aps.anl.gov/Accelerator-Operations-Physics/Software\#elegant}
  \item Manual:
        \url{https://ops.aps.anl.gov/manuals/elegant_latest/elegant.html}
  \item User forum:
        \url{https://www3.aps.anl.gov/forums/elegant/}
  \item \sdds{} information page:
        \url{https://ops.aps.anl.gov/SDDSInfo.shtml}
\end{itemize}

\section{Simulation Capabilities}

These capabilities correspond to commands in the main \elegant{} input file.

\begin{itemize}
  \item Tracking of rings, linacs, and beam-transport lines.
  \item Computation of $s$-dependent and final beam properties, including beam
        moments, Twiss parameters, matrix elements, the trajectory or closed
        orbit, the response matrix, and floor (survey) coordinates.
  \item Determination of apertures, including the physical (input) aperture,
        the dynamic acceptance, and the momentum acceptance.
  \item Optimization of lattice and beam properties.
  \item Tracking with errors, including generation and loading of error sets
        and correction of tunes, chromaticities, and orbit/trajectory.
  \item Control of lattice parameters, including scanning, altering, and
        loading values from \sdds{} data, as well as inserting and subdividing
        elements.
  \item Generation of particle bunches, or reading of an externally supplied
        distribution.
  \item Multi-stage simulation, in which results are passed from one run to the
        next.
  \item Time-dependent simulation through ramping and modulation of element
        parameters.
  \item Computation of frequency maps for rings.
  \item Changing of the particle type.
\end{itemize}

\section{Physics Capabilities}

These capabilities correspond to the elements available in the
lattice-definition file. They are grouped into single-particle dynamics,
collective dynamics, and other effects.

\subsection{Single-particle dynamics}
\begin{itemize}
  \item Magnetic elements, modeled with matrix and/or symplectic-integration
        methods.
  \item Stray fields and steering (corrector) elements.
  \item Undulator and wiggler kick maps.
  \item Synchrotron radiation, including its effect on the beam.
  \item RF cavities.
  \item Traveling-wave accelerating structures.
  \item Scattering from material and residual gas.
  \item Kickers.
  \item User-defined transformations and matrices.
  \item Apertures and scrapers.
\end{itemize}

\subsection{Collective dynamics}
\begin{itemize}
  \item Short-range wake fields and impedances.
  \item Long-range, resonant-mode wake fields.
  \item Coherent synchrotron radiation.
  \item Longitudinal space-charge impedance.
  \item Transverse space-charge kicks.
  \item Intrabeam scattering.
  \item Touschek scattering.
\end{itemize}

\subsection{Other}
\begin{itemize}
  \item Beam-position monitors.
  \item Phase-space output and analysis monitors.
  \item Digital feedback systems.
\end{itemize}

\section{Cooperative Capabilities}

These capabilities are provided by other programs that work together with
\elegant{} through \sdds{} files.

\begin{itemize}
  \item A general-purpose \sdds{} toolkit for postprocessing and graphics.
  \item Reading of particle data produced by ASTRA, IMPACT, and TRACK.
  \item Phase-space analysis.
  \item Computation of radiation properties, such as brightness and flux.
  \item Upsampling of particle distributions.
  \item Processing of quadrupole-scan emittance-measurement data.
  \item Computation of coherent-synchrotron-radiation impedance.
  \item Computation of lifetime, intrabeam-scattering, and potential-well
        effects.
  \item Translation of lattices to and from other formats.
  \item Computation of multipole-error data.
\end{itemize}

\nocite{SDDS1,SDDS2,Halbach_69a,Huang2004,Xiao2007SC,Xiao2008a,Xiao2010a,Warnock,Elleaume1992,ASTRA,IMPACT,TRACK}
\bibliographystyle{unsrt}
\bibliography{references}

\end{document}
